\section{Conclusion}\label{sec:nconclusion}

Sur les tabliers des notes de Grothendieck, les trois arches du \S\ref{sec:morearches} donnent des corollaires immédiats; nous les consignons parce que ce sont les premières entrées du registre tirées du fonds. La note du dossier 161-3, p.~26, contient deux exercices qui sont des ponts sur le tablier des préordres: le générateur retrouve le premier par l'arche de Morita, et sa forme relative aux préordres est démontrée en Lean; le second, sur les espaces, n'est pas formalisé. Les pseudo-corps sont le seul tablier du registre qui relie des théories équivalentes de façon non triviale. Nous en avons vérifié à la main la présentation par préfaisceaux, non celle par la topologie $\mathrm{ps.\,corps}$, et rien de ce qui les concerne n'est formalisé; l'argument par les points couvre les intersections d'une famille dénombrable de sous-topos de sa liste, non celles d'une famille non dénombrable. La profondeur d'un théorème engendré est bornée par celle de son arche et de son tablier \cite[transparent 8]{caramello-mines}: pour lire ces notes plus profondément, il faudrait des arches sur de vrais sites, que le registre ne contient pas.
